\documentclass[11pt]{article}
\usepackage[margin=0.85in]{geometry}
\usepackage{amsmath,amssymb,amsthm,longtable,booktabs}
\usepackage{algorithm,algpseudocode,float,hyperref}
\hypersetup{hidelinks}
\newtheorem{theorem}{Theorem}[section]
\newtheorem{lemma}[theorem]{Lemma}
\theoremstyle{definition}
\newtheorem{definition}[theorem]{Definition}
\newtheorem{remark}[theorem]{Remark}
\newcommand{\QQ}{\mathbb Q}
\newcommand{\RR}{\mathbb R}
\newcommand{\Eval}{\operatorname{IE}}
\newcommand{\width}{\operatorname{width}}
\newcommand{\Zero}{\operatorname{Zero}}
\newcommand{\prem}{\operatorname{prem}}
\newcommand{\res}{\operatorname{res}}
\newcommand{\lc}{\operatorname{lc}}
\newcommand{\disc}{\operatorname{disc}}
\newcommand{\alg}[1]{\textsc{#1}}
\title{Full Lazard CAD Using Rational Polynomials\\and Triangular Real-Point Certificates\\\large Executable pseudocode, correctness proofs, and audit requirements}
\author{CLPoly Project}
\date{October 22, 2026}
\begin{document}
\maketitle
\begin{abstract}
All stored polynomials in this specification have rational coefficients. A real
algebraic point is specified by a squarefree regular triangular chain and a
certified isolating box, following Xia--Zhang (2006). No instruction substitutes
an irrational number into a polynomial object. Specialization is represented by
keeping the rational polynomial together with the point certificate. We give
terminating algorithms for selected-point zero tests and signs, coefficientwise
zero tests, and Lazard evaluation, with proofs. Triangular decomposition routines
are imported explicitly from Xia--Yang (2016); isolation and refinement follow
Xia--Zhang (2006). Xia--Yang (2002) is supplementary only. The CAD implementation uses this rational-data representation. Complete lifting is specified and
proved relative to the explicitly imported projection theorem below.
These are mathematical proofs, not machine-checked verification of C++.
\end{abstract}
\tableofcontents

\section{Authority, scope, and permitted primitives}
This document specifies real-point operations, QQ Lazard evaluation, and complete
CAD lifting. The projection theorem is imported explicitly from the baseline
research and MPP; lifting, cell semantics and the driver are proved relative
to that theorem. Runtime polynomials always have rational coefficients.

The following cited algorithms are permitted subroutines; they are not unspecified
algebraic-number oracles. Their inputs, outputs, and restrictions are part of the
contract. All calls here have \emph{no free parameters} ($K=\QQ$, $u=\varnothing$).
\begin{longtable}{p{0.20\linewidth}p{0.32\linewidth}p{0.40\linewidth}}
\toprule Primitive & Exact source & Contract used here\\\midrule
\endhead
\alg{WRSD} & XY16 Algorithm 2.2, pp.40--41; Theorem 2.6 &
Returns $(H,G)$ with $V(T,c)=\bigcup_{C\in H}V(C)$ and
$V(T)\setminus V(c)=\bigcup_{C\in G}V(C)$.\\
\alg{RCtoSqrfree} & XY16 Algorithm 3.1 and Theorem 3.2, p.62 &
Finite squarefree regular decomposition of a zero-dimensional regular chain;
includes the first polynomial, not only higher levels.\\
\alg{NREALZERO} & XZ06 \S3, pp.855--858 &
Isolation of every real solution of a regular TSA, after the positive-root shifts
on p.856. Retain bound-polynomial and monotonicity certificates.\\
\alg{NSHR} & XZ06 pp.857--858 &
Refine a certified box from NREALZERO without changing its selected solution.
It is not an independent refinement of arbitrary coordinate intervals.\\
\alg{UNI} & Rational univariate isolation, as permitted in XZ06 \S3 &
All real roots of a nonzero squarefree rational polynomial, with positive
precision, in sorted strict brackets. Algorithm~\ref{a:uni} clears positive
denominators and supplies zero explicitly around the trusted nonzero-root isolator.\\
\bottomrule
\end{longtable}
We use polynomial arithmetic, rational comparisons, differentiation, coefficient
collection, and rational interval arithmetic as elementary primitives. No inverse
in a quotient algebra and no gcd over an algebraic coefficient field is needed.
XY02 REALZERO may report failure; it is not substituted for NREALZERO.
The legacy RSD implementation and its existing test remain available.
The new algorithms do not invoke it, directly or through their current
\texttt{charset.hh} call paths; RSD is not an imported primitive here.
The nonreduced-prefix case $T=[x^2,y^3]$, $c=xy$ prevents treating the
legacy backend as terminating on every regular input. The project note
\texttt{docs/cad/fixes/rsd-nonradical-prefix.md} records the counterexample
and the scope of the unresolved literature conditions.
The squarefree decomposition contract above remains a required interface.
The current C++ backend can repeat the same input on
$[x^2,y^3,z^2+xy]$: a valid WRSD zero branch need not simplify the prefix.
Until this progress issue is repaired, implementation conformance for arbitrary
regular-chain inputs is not established. This does not affect the proof below
that POINT-ZERO itself uses only the WRSD set contract.

\section{Representation and invariants}
Fix the \emph{lifting} coordinate order $x_1,\ldots,x_k$.
The public CAD argument is instead the \emph{projection} order. For a public
tuple $V=(v_1,\ldots,v_n)$, the local lifting notation means
$x_i=v_{n+1-i}$: projection eliminates $v_1,v_2,\ldots$, whereas lifting
first isolates $v_n$ and then extends towards $v_1$.
The C++ triangular-chain constructor receives the lifting tuple and a pointer
to an external lexicographic comparator. The tuple must be the reversed
contiguous suffix of the comparator's declared variable-priority list;
this also preserves the consecutive class indices used by WRSD. The main variable is
always the first, or smallest, variable under that comparator; the chain
equations and rational box remain in lifting order.
\begin{definition}[Valid real-point certificate]
A certificate $A=(T,B,\mathcal C)$ consists of rational polynomials
$T=[t_1(x_1),\ldots,t_k(x_1,\ldots,x_k)]$, a rational box
$B=\prod_i[a_i,b_i]$, and the isolation/refinement evidence $\mathcal C$.
Each $t_i$ has positive degree in $x_i$. The chain is zero-dimensional,
regular and squarefree: on \emph{every complex zero} of its preceding prefix,
$\lc(t_i,x_i)\disc(t_i,x_i)$ is nonzero; $t_1$ has no multiple root.
Every prefix box contains exactly one real zero of the corresponding prefix chain.
The selected point is in the interior of its box. Evidence $\mathcal C$ records
the rational coordinate shifts and the paired bound-root certificates used by
NREALZERO/NSHR. Rational roots can always be enclosed in strict rational brackets.
For $k=0$, $T$ and $B$ are empty and designate the unique empty point.
\end{definition}
An arbitrary box containing one solution is not automatically an NSHR certificate.
Constructor validation must establish the stated conditions or obtain a certificate
from NREALZERO. Rational coordinates may be represented by $q_i x_i-p_i=0$;
they are processed in exactly the same order as the other coordinates.
Each C++ triangular-chain object stores only a pointer to an external comparator.
Default copy and move operations preserve this pointer together with the
polynomial and monomial comparator pointers. Prefixes and decomposition branches
borrow the same comparator without rebinding or normalization. The caller must
keep it alive and its variable order unchanged for all chains, points and
converted polynomials that borrow it. The C++ extension routines
\texttt{append\_rational} and \texttt{roots\_over} receive this pointer
explicitly; no returned point borrows a comparator local to these routines.

A stored polynomial representative is always
\[p\in\QQ[x_1,\ldots,x_k,y].\]
For a prefix length $i$, define the \emph{semantic} map
\[E_i(p)=p(\alpha_1,\ldots,\alpha_i,x_{i+1},\ldots,x_k,y).
\]
$E_i$ is used in specifications and proofs only. It is never an executable
instruction. The representation $(p,A,i)$ has this meaning, while its data remains
rational. Variables $x_1,\ldots,x_i$ are frozen representative variables; the
remaining ones are active. A prefix certificate uses the restricted variable order,
not an order containing undeclared later variables.

\section{XZ06 isolation and refinement, made operational}
\subsection{Interval bounds and positive shifts}
For $q(u,z)=\sum_{j=0}^d q_j(u)z^j$ and a rational prefix box $X$, compute
$[c_j,d_j]=\Eval(q_j,X)$ and set
\[L(z)=\sum c_jz^j,\qquad U(z)=\sum d_jz^j.\]
XZ06 Proposition 3.1 gives $L^{(r)}(z)\le q^{(r)}(\alpha,z)\le U^{(r)}(z)$
for $z>0$ and all $r\ge0$. The same inequalities hold for all $u\in X$.
The bounds are rational polynomials, not necessarily integer polynomials.

To make the shift on XZ06 p.856 executable, first refine until
$0\notin\Eval(q_d,X)$. Let $m=\min(|c_d|,|d_d|)>0$ and
$M=\max_{j<d}\max(|c_j|,|d_j|)$. Choose rational
\[R>1+M/m,\qquad \widehat q(u,z)=q(u,z-R).
\]
Every real root in the old coordinate $y$ satisfies $|y|<R$, so its new
coordinate $z=y+R$ is strictly positive. This includes negative and zero roots.
The shift and its inverse must remain in the certificate. This is a concrete
root-bound choice for XZ06's shift, not a different isolation algorithm.
Indeed, for $|y|>1+M/m$, the magnitude of the leading term exceeds the sum of
all lower-term magnitudes by the geometric-series bound, so $q(u,y)\ne0$.

\begin{algorithm}[H]
\caption{XZ-LIFT: certified version of NREALZERO Step 2}\label{a:lift}
\begin{algorithmic}[1]
\Require Valid prefix certificate $A=(T,X,\mathcal C)$; $q$ with nonvanishing
leading coefficient and simple roots at the selected prefix
\Ensure All distinct roots of $q(\alpha,y)$, represented by $q,A$ and strict intervals
\State Determine and record the shift $R$ above; set $\widehat q(u,z)=q(u,z-R)$
\State $n\gets0$
\Loop
 \State $(L,U)\gets$ coefficient interval bounds of $\widehat q$ on $X$
 \State $n\gets n+1$
 \If{$L,U$ do not have degree $d$, equal leading signs, or are not squarefree}
  \State $X\gets\alg{NSHR}(A)$; \textbf{continue}
 \EndIf
 \State $S_L\gets\alg{UNI}(L,2^{-n})$; $S_U\gets\alg{UNI}(U,2^{-n})$
 \If{either list has a nonpositive root or their lengths differ}
  \State $X\gets\alg{NSHR}(A)$; \textbf{continue}
 \EndIf
 \State Pair roots by increasing order; let $J_j$ be the hull of their strict brackets
 \If{the $J_j$ are positive and disjoint, and $L',U'$ have the same strict sign on each $J_j$}
  \State \Return $(q,A,\{J_j-R\}_j)$ and bound/shift certificates
 \EndIf
 \State $X\gets\alg{NSHR}(A)$
\EndLoop
\end{algorithmic}
\end{algorithm}
The derivative acceptance tests use rational interval evaluation of $L'$ and $U'$.
For empty prefixes NSHR is a no-op; univariate bracket precision still increases.
A constant nonzero $q$ returns the empty list before calling this algorithm. A zero
specialization is a separate error, never an empty finite root set. The additional
squarefree and strict-sign checks are explicit sufficient certificates for XZ06's
pairing/monotonicity conditions; they are not claimed to be verbatim paper lines.

\begin{lemma}[XZ-LIFT correctness and termination]\label{l:lift}
Algorithm~\ref{a:lift} returns exactly all roots, without requiring irrational
coefficient operations.
\end{lemma}
\begin{proof}
At acceptance, both bounding polynomials have the same number of simple real
roots, all positive, the same degree and leading sign. Outside the paired hulls,
they have the same nonzero sign: count crossings from $+\infty$, one simple root
from each polynomial per hull. Thus the target polynomial cannot vanish outside
these hulls. On each hull the common strict derivative sign implies strict
monotonicity of the target, by XZ06 Proposition 3.1. At the two endpoints both
bounds have opposite signs, hence so does the target. The intermediate value
theorem gives exactly one root. Shifting back is a bijection.

As the prefix shrinks to its selected point, coefficient interval widths tend
to zero. Both bound polynomials converge to the same fixed polynomial of degree
$d$ with simple roots, none at zero. In a sufficiently small coefficient
neighborhood, each simple real root persists and stays separated, and each
nonreal root remains nonreal. Their ordered real roots therefore pair and converge.
The univariate bracket widths also tend to zero. Derivatives at the limiting roots
are nonzero, so the strict interval derivative certificates eventually hold.
All acceptance conditions are eventually met. This is the stability argument
underlying XZ06 \S3, with explicit sufficient acceptance predicates.
\end{proof}

\subsection{NSHR and arbitrary requested precision}
For review, the imported NSHR steps are: shrink the first-coordinate bracket to
at most half its width by rational univariate bracketing; for each subsequent
coordinate, compute new lower and upper bound polynomials using the already
refined prefix; bracket their respective unique roots inside the old coordinate
interval to width at most $1/8$ of that interval; take the hull of these two
brackets. Retain the same shift and equation. The preceding strict derivative
certificates and monotonicity of interval evaluation preserve the required roots
and monotonicity on every later call.

Univariate bracketing is finite rational bisection of a unique simple root with
opposite endpoint signs. If the midpoint is exactly a root, keep strict endpoints
on its two sides rather than returning a degenerate bracket. Such endpoints may
be chosen arbitrarily close to the rational midpoint; the established
monotonicity gives their signs. Thus no point interval is required.

\begin{algorithm}[H]
\caption{NSHR-QQ: explicit rational-data transcription of XZ06 NSHR}\label{a:nshr}
\begin{algorithmic}[1]
\Require Valid certificate $A=(T,B,\mathcal C)$ with its per-level bound/shift evidence
\Ensure Nested valid box selecting the same point
\If{$T=[]$} \State \Return $A$ \EndIf
\State Save old box $B^{old}$
\State $B_1\gets\alg{BRACKET}(t_1,B^{old}_1,\width(B^{old}_1)/2)$
\For{$j=2,\ldots,k$}
 \State $R_j\gets$ recorded shift for level $j$
 \State $q(u,z)\gets t_j(u,z-R_j)$ by rational polynomial translation
 \State $(L,U)\gets$ coefficient interval bounds of $q$ on the updated prefix $B_{<j}$
 \State $J\gets B^{old}_j+R_j$; $\epsilon\gets\width(B^{old}_j)/8$
 \State $I_L\gets\alg{BRACKET}(L,J,\epsilon)$; $I_U\gets\alg{BRACKET}(U,J,\epsilon)$
 \State $B_j\gets\operatorname{hull}(I_L,I_U)-R_j$; update bound evidence
\EndFor
\State \Return updated $A$
\end{algorithmic}
\end{algorithm}
\alg{BRACKET} is the strict rational bisection procedure described above: while
width exceeds the bound, evaluate the midpoint rationally and retain the
opposite-sign half; for an exact midpoint root choose strict brackets on its two
sides of width at most half the old width and at most the requested width.
The input polynomial is strictly monotone on the input interval and has opposite
endpoint signs, so every step is defined. For level one a univariate root
certificate supplies the same unique-root bracketing property; monotonicity can
be certified during construction by refining its simple-root interval.

\begin{algorithm}[H]
\caption{REFINE: precision wrapper around XZ06 NSHR}\label{a:refine}
\begin{algorithmic}[1]
\Require Valid $A$, rational $\epsilon>0$
\Ensure Same selected point, nested certified box, every width $<\epsilon$
\While{$\max_j\width(B_j)\ge\epsilon$}
 \State $A\gets\alg{NSHR}(A)$ \Comment{XZ06 pp.857--858; empty point needs no loop}
\EndWhile
\State \Return $A$
\end{algorithmic}
\end{algorithm}
\begin{lemma}[REFINE terminates]\label{l:refine}
Successive NSHR calls preserve the point and make every coordinate width tend to zero.
\end{lemma}
\begin{proof}
Point preservation is XZ06's NSHR argument: the updated bounds lie between the
old bounds, and their unique roots bracket the target root. Their hull is inside
the old interval. The first width decreases by a factor at most $1/2$.
Inductively suppose all earlier widths tend to zero. The two bounding roots at
level $j$ converge to the same simple target root, so their distance $\delta_n$
tends to zero. Since each bracket width is at most $w_{j,n}/8$, the next hull
satisfies $w_{j,n+1}\le\delta_n+w_{j,n}/4$. These nonnegative widths are nested;
taking limits gives $w_j\le w_j/4$, hence $w_j=0$. Induction proves the claim.
Do not replace this argument by the false assertion that every NSHR call halves
\emph{every} coordinate width.
\end{proof}

\begin{algorithm}[H]
\caption{ISOLATE-REGULAR: equation-only NREALZERO driver}\label{a:isolate}
\begin{algorithmic}[1]
\Require A squarefree zero-dimensional regular rational chain $T$ in the stated order
\Ensure Certificates for all distinct real solutions of $T$
\If{$T=[]$} \State \Return the singleton list containing the empty point \EndIf
\State $L\gets\alg{UNI}(t_1)$; refine its brackets until $0\notin\Eval(t'_1,J)$
\State Attach $[t_1]$ and the strict root evidence to each interval in $L$
\For{$j=2,\ldots,k$}
 \State $L_{new}\gets[]$
 \For{each prefix certificate $A$ in $L$}
  \State $S\gets\alg{XZ-LIFT}(t_j,A)$
  \State Append every extended certificate $(A,t_j,J)$, $J\in S$, to $L_{new}$
 \EndFor
 \State $L\gets L_{new}$
\EndFor
\State \Return $L$
\end{algorithmic}
\end{algorithm}
This is NREALZERO for the pure-equation case needed here. Inequality filtering
(Step 3 in the paper) is not invoked. The first coordinate is isolated on the full
real line by UNI; later levels use the recorded positive shifts. The free functions \texttt{shrink} and \texttt{refine} implement Algorithms~\ref{a:nshr} and~\ref{a:refine};
the point member functions only forward to them. The NSHR body
works in original coordinates with a shift only for the active bound polynomial.
By induction on $j$, Lemma~\ref{l:lift} gives all extensions of each prefix and no
others. Distinct prefixes or distinct intervals within a fibre give distinct points.
The regular squarefree hypothesis establishes every XZ-LIFT precondition.

\section{Selected-point zero and sign without extension arithmetic}
WRSD directly returns two families $(Z,N)$ with
\[\bigcup_{C\in Z}V(C)=V(T)\cap V(c),\qquad
  \bigcup_{C\in N}V(C)=V(T)\setminus V(c).\]
These set equalities are the entire decomposition contract used below.
The zero label does not require $\prem(c;C)=0$ and is never reconstructed
from a pseudo-remainder test.
Branches need not be pairwise disjoint and need not themselves be squarefree:
we refine only the original valid certificate, never a newly returned branch box.

\begin{algorithm}[H]
\caption{POINT-ZERO: exact zero at one selected branch}\label{a:zero}
\begin{algorithmic}[1]
\Require Valid $A=(T,B,\mathcal C)$ of dimension $i$;
$c\in\QQ[x_1,\ldots,x_i]$
\Ensure Returns true exactly when $c(\alpha)=0$; may refine $B$
\If{$i=0$} \State \Return $(c=0)$ by rational comparison \EndIf
\State $(Z,N)\gets\alg{WRSD}(T,c)$
\Loop
 \State Remove from $Z$ and $N$ each $C$ with some $g\in C$ satisfying $0\notin\Eval(g,B)$
 \State require $Z\cup N\ne\varnothing$
 \If{$N=\varnothing$} \State \Return true \EndIf
 \If{$Z=\varnothing$} \State \Return false \EndIf
 \State $A\gets\alg{REFINE}(A,\frac12\max_j\width(B_j))$
\EndLoop
\end{algorithmic}
\end{algorithm}
All equations in $C$ are evaluated on the original coordinate box, not on separate
univariate resultant-root candidates. A component is removed only by a certified
nonzero equation; interval overlap alone never proves membership.

\begin{theorem}[POINT-ZERO is a decision procedure]\label{t:zero}
Algorithm~\ref{a:zero} terminates and its Boolean result is correct.
\end{theorem}
\begin{proof}
WRSD terminates and supplies finitely many regular components with the two
set equalities stated above. A zero-labelled component consists entirely of
zeros of $c$, and a nonzero-labelled component has no zero of $c$.
Thus no point lies in components with opposite labels. No assertion about
membership in a nonradical saturated ideal is needed.

At least one component contains $\alpha$. None of its equations can have an
interval enclosure excluding zero, so it is never removed. Every component not
containing $\alpha$ has an equation $g$ with $g(\alpha)\ne0$. Polynomial interval
evaluation is inclusion preserving and converges on shrinking boxes; hence that
component is eventually removed. There are finitely many components. Eventually
all remaining components contain $\alpha$ and have the same label, forcing the
correct return. REFINE terminates by Lemma~\ref{l:refine}.
\end{proof}

\begin{algorithm}[H]
\caption{POINT-SIGN}\label{a:sign}
\begin{algorithmic}[1]
\Require Valid $A$; rational polynomial $c$ in its coordinates
\Ensure Exact sign of $c(\alpha)$
\State $[l,u]\gets\Eval(c,B)$
\If{$l>0$} \State \Return $+1$ \EndIf
\If{$u<0$} \State \Return $-1$ \EndIf
\If{$\alg{POINT-ZERO}(c,A)$} \State \Return $0$ \EndIf
\Loop
 \State $[l,u]\gets\Eval(c,B)$
 \If{$l>0$} \State \Return $+1$ \EndIf
 \If{$u<0$} \State \Return $-1$ \EndIf
 \State $A\gets\alg{REFINE}(A,\frac12\max_j\width(B_j))$
\EndLoop
\end{algorithmic}
\end{algorithm}
For dimension zero the interval is the exact rational singleton and a nonzero
constant returns before refinement. Correctness follows from inclusion;
the initial strict-sign path is immediate, and the remaining path terminates
by the exact nonzero decision and interval convergence. An interval such as $[0,b]$, $b>0$, proves neither zero nor a strict sign.

\section{Executable coefficientwise specialization test}
\begin{algorithm}[H]
\caption{PREFIX-ZERO: no irrational substitution}\label{a:prefix}
\begin{algorithmic}[1]
\Require $p\in\QQ[x_1,\ldots,x_k,y]$; valid prefix $A_i$; active variables
$W=(x_{i+1},\ldots,x_k,y)$
\Ensure Returns true exactly when $E_i(p)$ is the zero polynomial
\State $M\gets$ empty coefficient map (missing entries mean zero)
\For{each term $a x^\beta W^\gamma$ of $p$}
 \State $M[\gamma]\gets M[\gamma]+a x^\beta$ \Comment{collect equal active monomials first}
\EndFor
\For{each nonzero collected coefficient $c$ in $M$}
 \If{not $\alg{POINT-ZERO}(c,A_i)$} \State \Return false \EndIf
\EndFor
\State \Return true
\end{algorithmic}
\end{algorithm}
\begin{lemma}
PREFIX-ZERO terminates and has its stated meaning.
\end{lemma}
\begin{proof}
Collection is a finite rational polynomial computation and yields
$p=\sum_\gamma c_\gamma(x_1,\ldots,x_i)W^\gamma$.
Distinct monomials in the still formal variables $W$ are linearly independent.
Consequently $E_i(p)=0$ if and only if every $c_\gamma(\alpha)=0$.
Apply Theorem~\ref{t:zero} to the finitely many coefficients.
\end{proof}

\section{Lazard definition, full pseudocode, and proof}
For a nonzero rational polynomial $f$, define its Taylor expansion semantically by
\[
 f=\sum_{\beta\in\mathbb N^k}c_\beta(y)
       \prod_{j=1}^k(x_j-\alpha_j)^{\beta_j}.
\]
The support $S=\{\beta:c_\beta\ne0\}$ is finite and nonempty because translation
is invertible. The valuation is $v=\min_{\rm lex}S$ with $x_1$ highest priority;
the Lazard value is $f_\alpha=c_v(y)$. This defines the result; no Taylor expansion
with irrational coefficients is constructed.

\begin{algorithm}[H]
\caption{LAZARD-QQ: frozen rational representatives}\label{a:lazard}
\begin{algorithmic}[1]
\Require Nonzero $f\in\QQ[x_1,\ldots,x_k,y]$; valid $A$; this explicit coordinate order
\Ensure Return $(p,v)$, interpreted relative to the unchanged input point $A$, with
$E_k(p)=f_\alpha\ne0$ and $v=v_\alpha(f)$
\State $p\gets f$; $v\gets[]$
\For{$i=1,\ldots,k$}
 \State $A_i\gets$ certified prefix of $A$ in its restricted variable order
 \State $h\gets p$; $r\gets0$; $b\gets1$
 \While{$\alg{PREFIX-ZERO}(h,A_i,(x_{i+1},\ldots,x_k,y))$}
  \State $h\gets\partial h/\partial x_i$ \Comment{differentiate the unreduced rational representative}
  \State $r\gets r+1$; $b\gets b r$
 \EndWhile
 \State $p\gets h/b$; append $r$ to $v$ \Comment{division only by a nonzero integer}
\EndFor
\State \Return $(p,v)$
\end{algorithmic}
\end{algorithm}
Reject zero input before line 1. For $k=0$ return $(f,[])$.
The caller retains $A$; the returned C++ object owns only its QQ representative,
variable comparator and valuation. It is not a self-contained algebraic value,
and must not be reinterpreted using a different conjugate point.
Prefix refinements may be kept as local copies or propagated to the parent box;
coordinate nesting preserves the point either way. No equation of $T$ is used to
reduce $p$ or $h$. In particular, the algorithm neither divides by
$x_i-\alpha_i$ nor drops a pseudo-division multiplier. Optional future reductions
require a separate proof that they commute with all remaining derivatives.

\begin{lemma}[Derivative extraction]
For a nonzero $g(z)=\sum_{j=0}^d a_j(z-\theta)^j$ over a polynomial ring with real
coefficients, the first nonzero derivative specialized at $z=\theta$ is
$g^{(r)}(\theta)=r!a_r$, where $r$ is the first nonzero coefficient index.
\end{lemma}
\begin{proof}
Differentiating $r$ times and setting $z=\theta$ kills every translated monomial
except exponent $r$, whose multiplier is $r!$. Some coefficient is nonzero and
factorials are invertible in the real coefficient ring used here.
\end{proof}

\begin{theorem}[LAZARD-QQ total correctness]\label{t:lazard}
Algorithm~\ref{a:lazard} terminates and meets its postcondition. Every computed
polynomial has rational coefficients.
\end{theorem}
\begin{proof}
After iteration $i$ the invariant is
\[
 E_i(p)=\sum_{\substack{\beta\in S\\\beta_1=v_1,\ldots,\beta_i=v_i}}
 c_\beta(y)\prod_{j>i}(x_j-\alpha_j)^{\beta_j}\ne0.
\]
It holds initially with the entire Taylor expansion. At iteration $i$, earlier
specialization commutes with differentiation in the distinct variable $x_i$:
$E_{i-1}(\partial_{x_i}^r p)=\partial_{x_i}^r E_{i-1}(p)$.
PREFIX-ZERO tests exactly whether the next specialization of this derivative is
zero as a polynomial in the later variables. Derivative extraction therefore
selects the smallest possible exponent $v_i$ and division by $v_i!$ gives the
asserted invariant. Keeping the earlier variables in the representative does not
change this semantic identity.

Successive smallest-prefix choices are exactly lexicographic minimization. At
$i=k$ the invariant is $E_k(p)=c_v(y)\ne0$. Each inner loop has at most
$\deg_{x_i}(p)$ differentiations, since the nonzero semantic polynomial has no
larger degree than its representative. Every zero test terminates. The outer loop
has $k$ steps, so the algorithm terminates. Differentiation, coefficient grouping,
and division by factorials preserve rational coefficients. All interval and
triangular subroutines also operate on rational data.
\end{proof}

\section{Output use in lifting: branch restriction before root isolation}
Lifting interprets each QQ representative at the selected prefix point. Branch
restriction preserves the zero-dimensional hypotheses required by the cited
triangular isolation algorithms.

A point belongs to a candidate component $C$ exactly when POINT-ZERO returns
true for every $g\in C$. This terminates and works even when components overlap.
The selected prefix is re-encoded recursively, then only its next coordinate
is lifted. No interval-overlap membership test is used.

Write $\alg{LIFT-ONE}(C,P)$ for one execution of NREALZERO Step 2: $C$ is a
squarefree regular chain of length $k$, and $P$ is a valid certificate on its
length-$k-1$ prefix. Call XZ-LIFT on the last equation at $P$, and append each
returned strict root bracket and its fixed shift to the refined prefix.
The inherited prefix certificates and the new XZ-LIFT evidence certify the
combined point, so no already verified prefix equation needs to be re-certified.
This helper is implemented by \texttt{\_\_lift\_real\_points}; full-chain
ISOLATE-REGULAR calls it on every prefix, while ROOTS-QQ calls it only on its
selected prefix.

\begin{algorithm}[H]
\caption{REENCODE: retain the selected point on another regular chain}\label{a:reencode}
\begin{algorithmic}[1]
\Require Valid $A=(T,B,\mathcal C)$; regular chain $C$ on the same $k$ coordinates
\Ensure For $\alpha\in V(C)$, a squarefree certificate for $\alpha$; otherwise reject
\If{$C=T$ as normalized equation lists} \State \Return $A$ \EndIf
\State Require POINT-ZERO to hold on $A$ for every equation of $C$
\If{$k=0$} \State \Return $A$ \EndIf
\For{$D\in\alg{RCtoSqrfree}(C)$}
 \If{$C$ was not squarefree and some equation of $D$ fails POINT-ZERO at $A$}
  \State \textbf{continue}
 \EndIf
 \State $P\gets\alg{REENCODE}(A_{<k},D_{<k})$
 \For{$Q\in\alg{LIFT-ONE}(D,P)$}
  \If{$\alg{SAME-POINT}(Q,A)$} \State \Return $Q$ \EndIf
 \EndFor
\EndFor
\State Report loss of the selected point
\end{algorithmic}
\end{algorithm}
SAME-POINT first tests every equation of the second chain at the first point
using POINT-ZERO. If any test fails, the points differ. Otherwise refine the
first certificate until its box is contained in the interior of the second box
(equal) or is disjoint from it (different). The first point then lies on the
second chain, whose certificate has one interior zero and none on its boundary.
This proves correctness and termination of the comparison, without assuming
that overlapping boxes represent the same point.

\begin{lemma}[Selected-prefix re-encoding]
REENCODE terminates and returns the same point on a squarefree component of $C$.
\end{lemma}
\begin{proof}
Induct on $k$. Identical normalized chains preserve the existing certificate;
$k=0$ is the empty point. Otherwise the initial membership tests verify the
precondition, and RCtoSqrfree supplies finitely many covering squarefree chains.
At least one contains $\alpha$; the branch tests remove only those that do not.
By induction, its recursive call gives exactly the first $k-1$ coordinates of
$\alpha$. XZ-LIFT then enumerates all last-coordinate roots on this prefix,
including $\alpha_k$. SAME-POINT accepts exactly the desired full point.
All branch lists and root lists are finite, and recursion lowers $k$.
No global containment $V(C)\subseteq V(T)$ is needed because candidate membership
is tested before box matching. Independent chain/point prefix operations only
restrict already verified certificates and preserve the external comparator pointer.
\end{proof}

\begin{algorithm}[H]
\caption{ROOTS-QQ: isolate roots only over the selected prefix}\label{a:roots}
\begin{algorithmic}[1]
\Require Valid $A=(T,B,\mathcal C)$; $p=\sum_{j=0}^d c_j(x)y^j$
\Ensure For $E_k(p)\ne0$, all distinct fibre roots; reject an identically zero fibre
\State $C\gets T$; $e\gets-1$
\For{$j=d,d-1,\ldots,0$}
 \State $(H,G)\gets\alg{WRSD}(C,c_j)$
 \State Choose $D\in H\cup G$ containing $\alpha$ by POINT-ZERO on its equations
 \State $C\gets D$
 \If{$D\in G$} \State $e\gets j$; \textbf{break} \EndIf
\EndFor
\If{$e<0$} \State Reject the identically zero specialization \EndIf
\If{$e=0$} \State \Return empty list \EndIf
\State $q\gets\sum_{j=0}^e c_jy^j$; $R\gets[]$
\For{$D\in\alg{RCtoSqrfree}([C,q])$}
 \If{some equation of $D_{\le k}$ fails POINT-ZERO at $A$} \State \textbf{continue} \EndIf
 \State $P\gets\alg{REENCODE}(A,D_{\le k})$
 \For{$Q\in\alg{LIFT-ONE}(D,P)$}
  \State $duplicate\gets\mathrm{false}$
  \For{$Q'\in R$}
   \If{their last-coordinate brackets are disjoint} \State \textbf{continue} \EndIf
   \If{$\alg{SAME-POINT}(Q,Q')$}
    \State $duplicate\gets\mathrm{true}$; \textbf{break}
   \EndIf
   \While{their last-coordinate brackets overlap}
    \State Refine $Q$ and $Q'$ to half their respective maximum widths
   \EndWhile
  \EndFor
  \If{not $duplicate$} \State Append $Q$ to $R$ \EndIf
 \EndFor
\EndFor
\State Sort $R$ by the left endpoints of its disjoint last brackets; \Return $R$
\end{algorithmic}
\end{algorithm}
For an empty parent, coefficient labels are rational comparisons; no triangular
routine is called on an empty chain. LIFT-ONE uses XZ-LIFT with the empty prefix
and the rational UNI adapter, so this case is included in the same control flow.

\begin{lemma}[ROOTS-QQ contract]
All calls in Algorithm~\ref{a:roots} satisfy their algebraic preconditions, and
its result is complete and exact.
\end{lemma}
\begin{proof}
WRSD covers the current chain and partitions its zeros by each coefficient's
vanishing. The selected point persists. On the selected chain all skipped higher
coefficients vanish at every zero, and the first nonzero-labelled $c_e$ is
nonzero at every zero. If no such coefficient exists the fibre is identically
zero; if $e=0$ it has no roots. Otherwise $[C,q]$ is zero-dimensional and regular.
RCtoSqrfree preserves its zeros and ensures all LIFT-ONE simple-root hypotheses.
A branch whose prefix does not contain $\alpha$ contributes no target fibre
root and can be discarded before isolation. On each retained branch, REENCODE
preserves the selected prefix and XZ-LIFT returns all last-coordinate roots on
that prefix, so the union is exactly the target root set. There is no need to
isolate real points over other prefixes.

All candidates have the same first $k$ coordinates. Disjoint last brackets
therefore prove inequality; overlapping ones use the exact SAME-POINT test.
Only equal points are removed. For distinct last coordinates, convergence and
positive separation make the persistent refinements terminate. After each
insertion all stored last brackets are pairwise disjoint; later refinements are
nested and preserve that invariant. Sorting by their left endpoints is a strict
order comparison with no algebraic work inside the comparator. Each list and
loop is finite, and every invoked decision/refinement terminates.
\end{proof}
This construction is implemented by \texttt{roots\_over}; its correctness is
relative to the same XY16 and XZ06 contracts as full-chain isolation. Taking
$T\cup\{p\}$ without first restricting the degree branch can produce a
positive-dimensional system and is forbidden.

\section{UNI-QQ: the checked rational univariate adapter}
\label{s:uni}
The scalar integer \texttt{uspensky} interface inherited from the trusted
baseline returns the \emph{nonzero} real roots of a squarefree polynomial.
It is not by itself the UNI contract used here. The following adapter supplies
zero explicitly. UNI in Algorithms 1--10 means this checked adapter; its input
must be squarefree and its requested precision positive. A nonzero constant is
allowed. Zero and nonsquarefree inputs are rejected, not silently interpreted as
finite empty root sets. All XZ-LIFT callers check squarefreeness before UNI, and
ISOLATE-REGULAR has a squarefree chain as its precondition.

\begin{algorithm}[H]
\caption{UNI-QQ: all simple real roots with strict brackets}\label{a:uni}
\begin{algorithmic}[1]
\Require Rational coefficient vector $f$; rational precision $\epsilon$
\Ensure For squarefree nonzero $f$ and $\epsilon>0$, all real roots in sorted,
strictly separated brackets of positive width at most $\epsilon$, with strict derivative sign
\State Reject $\epsilon\le0$; clear a common positive denominator to obtain $z\in\mathbb Z[t]$
\State Reject $z=0$; if $\deg z=0$, return $[]$
\State Reject unless $\deg\gcd(z,z')=0$
\State $L\gets\alg{ScalarUspensky}(z)$ \Comment{trusted nonzero-root interface}
\If{$f(0)=0$} \State append $[0,0]$ to $L$ \EndIf
\For{each interval $J\in L$}
 \If{an endpoint $r$ of $J$ is an exact root}
  \State $\delta\gets\epsilon/4$
  \Repeat
   \State $J\gets[r-\delta,r+\delta]$; $\delta\gets\delta/2$
  \Until{$\Eval(f',J)$ has a strict sign and $f(\inf J)f(\sup J)<0$}
 \EndIf
 \State $J\gets\alg{BRACKET}(f,J,\epsilon)$
 \While{$\Eval(f',J)$ has no strict sign}
  \State $J\gets\alg{BRACKET}(f,J,\width(J)/2)$
 \EndWhile
\EndFor
\State Sort $L$ by left endpoints
\For{each adjacent pair $J,K$ in $L$}
 \While{$J\cap K\ne\varnothing$}
  \State Replace each by its BRACKET refinement to half its width
 \EndWhile
\EndFor
\State \Return $L$
\end{algorithmic}
\end{algorithm}

\begin{lemma}[UNI adapter correctness and termination]\label{l:uni}
Algorithm~\ref{a:uni} satisfies its checked contract.
\end{lemma}
\begin{proof}
Positive denominator clearing preserves the roots. The gcd test ensures every
root is simple. ScalarUspensky supplies each nonzero root exactly once, so the
additional zero singleton supplies precisely the missing root, without duplicates.
At every simple root the derivative is nonzero; hence sufficiently small rational
neighborhoods have strict derivative sign and opposite endpoint signs. This
proves termination and validity of the endpoint conversion. A strict isolating
bracket of a simple root can be bisected even before a derivative sign has been
certified: the unique root is retained by opposite endpoint signs. If the
midpoint equals that root, the chosen smaller symmetric bracket lies inside the
old one and contains no other root. Subsequent bisection makes the derivative
interval converge to its nonzero value. Distinct roots have positive separation,
so neighboring brackets eventually become disjoint. Later refinements are nested
and cannot destroy earlier separations. Every operation is rational and finite.
\end{proof}

\section{Complete CAD: projection contract and cell semantics}
\label{s:cad}
This section closes the lifting proof using the projection/delineability theorem
already referenced in the baseline project research. It does not replace XY16 or
XZ06 with another algebraic-number algorithm. We explicitly import the mathematical
result of McCallum--Parusi\'nski--Paunescu (MPP)\cite{MPP}, Theorem 5.1 and
\S5.3, and the trusted baseline projection, factorization and tree-storage contracts.
The following is an implementation-level refinement of MPP Algorithm 2.12.
It is a mathematical proof relative to these imports, not a machine proof of them.

\subsection{Projection data and the imported theorem}
Write $x_1,\ldots,x_n$ for the \emph{lifting} order; the public API receives the
reverse \emph{projection} order $(x_n,\ldots,x_1)$. Set $A_n$ to the input after
clearing positive rational denominators, retaining the original rational inputs
separately. At stage $k$, the trusted FACTOR projection supplies $B_k$, the
pairwise nonassociate primitive irreducible factors of positive degree in $x_k$,
and passes to $A_{k-1}$ the contents/lower-variable factors and the nonconstant
leading coefficients, trailing coefficients, discriminants and pairwise
resultants of $B_k$. Zero input polynomials are discarded for projection and
retained for final signs; nonzero constants need no projection.
Thus every nonzero $a\in A_k$ has a factorization
\[
 a=c_a\prod_{b\in B_k}b^{e_{a,b}},\qquad e_{a,b}\ge0,
\]
where $c_a$ depends only on $x_1,\ldots,x_{k-1}$ and is a nonzero constant times
products of data passed to the lower-dimensional problem. Signs/units of these
factorizations need not be stored because final signs use the original inputs.
The baseline \texttt{\_\_project\_full(...,LAZARD,FACTOR)} provides this contract;
\texttt{cad\_tree} reverses the layer array to give $B_1,\ldots,B_n$.

MPP Theorem 5.1 states that a polynomial with positive last-variable degree and
nonzero discriminant, leading coefficient and constant-in-the-last-variable
coefficient is Lazard analytically delineable on a connected analytic submanifold
where those three polynomials are valuation-invariant. Its moreover clause covers
multiplication by $x_k$. Here ``nonzero'' means nonzero as a polynomial, not
nonvanishing at every point of the cell. Nullifying lower-dimensional cells are
therefore not excluded. MPP Definition 2.10 and Proposition 2.11 identify the
continuous ordered root functions, their constant multiplicities, and valuation
invariance on their sections and sectors.

To apply the theorem to $B_k$, let $H$ be the product of its factors other than
the possible associate of $x_k$. Irreducibility ensures that every factor of $H$
has nonzero trailing coefficient, and squarefreeness and pairwise coprimality
ensure that $\disc(H)$ is a nonzero polynomial. Its leading and trailing
coefficients are products of those of the factors, while
\[
 \disc(H)=\pm\prod_b\disc(b)\prod_{b<c}\res(b,c)^2.
\]
Consequently the projection hypotheses imply valuation invariance of these
three polynomials on the parent cell. MPP Theorem 5.1 applies to $H$ and, if
needed, to $x_kH$. This is the product argument in MPP \S5.3. If $B_k$ is empty,
there are no sections; if it consists only of an associate of $x_k$, its sole
section is $x_k=0$ and the claim is immediate. Distinct factors may have coincident
specialized roots on lower-dimensional parents; we always use the distinct roots
of their product, with no assumption that individual specialized lists are disjoint.

\subsection{The meaning of stored nodes}
The unique depth-zero cell is $\mathbb R^0$, with the empty point certificate.
Suppose a parent node denotes a connected analytic semialgebraic cell $S$ with
sample $\alpha\in S$. The product theorem gives a fixed number $m$ of ordered
analytic functions $\theta_1<\cdots<\theta_m:S\to\mathbb R$ whose values at every
$u\in S$ are the distinct roots of the product of the Lazard evaluations of $B_k$.
The $j$th section denotes
\[
 \Gamma_j=\{(u,\theta_j(u)):u\in S\},
\]
and the sectors denote
\[
 \Omega_j=\{(u,t):u\in S,\ \theta_j(u)<t<\theta_{j+1}(u)\},
 \quad 0\le j\le m,
\]
where $\theta_0=-\infty$ and $\theta_{m+1}=+\infty$.
With $m=0$ this is the single cylinder $S\times\mathbb R$.
The runtime root certificates describe only $(\alpha,\theta_j(\alpha))$; they
are not asserted to be defining equations of the entire section.

The runtime root table is ordered by the distinct roots of the product.
Stored boundaries retain the baseline \texttt{cad\_root(b,r)} meaning: source
polynomial index $b$ in the current layer, and zero-based distinct-root index $r$
within that source's Lazard evaluation at this parent. For the $j$th product
root, set $sources_j$ to its nonempty source set and let
\[
 \nu_b(j)=\#\{h<j:b\in sources_h\}.
\]
Choose the least $b\in sources_j$ for the stored reference $(b,\nu_b(j))$.
Shared roots advance the counters of every source, so later roots of each
source retain their correct ordinals. The \texttt{root} query, given the level, parent and reference, resolves this identity using the source
lists, including valid alternative source identities for a shared root.
The original \texttt{cad\_tree} and its boundary type are not reinterpreted as
local-union indices. The parent, source and per-source ordinal identify the
same continuous section function guaranteed by projection, while the parent
stack retains the merged sample root list.

These cells are semialgebraic: the prefix valuation $r$ of the product is constant
on $S$, so its Lazard root relation is the zero set, over $S$, of the rational
polynomial $(\partial_1^{r_1}\cdots\partial_{k-1}^{r_{k-1}}H)/\prod r_i!$
(or the corresponding product including $x_k$). Choosing its $j$th distinct real
root is a first-order real-algebraic condition. Closure of semialgebraic sets under
projection yields the graph; this standard closure theorem is an explicit
mathematical import, not an executable quantifier-elimination step.
The sectors are semialgebraic by the same closure property.

\section{Complete CAD: executable lifting and driver}
\subsection{Simplest rational samples in open intervals}\label{s:sampling}
We order reduced rationals first by increasing positive denominator and then by
increasing absolute numerator, as in SageMath's \texttt{simplest\_rational}
specification. The following exact continued-fraction construction uses the
integer-part/reciprocal method also described by \texttt{approxRational} in
Haskell's \texttt{Data.Ratio} module, but with strictly open endpoints.
These are implementation references, not new trusted mathematical oracles:
\href{https://doc.sagemath.org/html/en/reference/rings_numerical/sage/rings/real_mpfi.html}{SageMath interval documentation} and
\href{https://hackage-content.haskell.org/package/base-4.7.0.1/docs/src/Data-Ratio.html}{Data.Ratio source}.
No floating-point operations are required. All integer parts use exact floor
division of the QQ numerator by its positive denominator.

\begin{samepage}
\paragraph{Auxiliary procedure SAMPLE-RATIONAL.}
Inputs are open rational endpoints; a missing left or right endpoint denotes
$-\infty$ or $+\infty$, respectively. Finite equal or reversed endpoints are
rejected. Its implementation is \texttt{sample\_rational} in the independent
\texttt{rational\_sample.hh} header.
\begin{algorithmic}[1]
\If{$a,b$ are both finite}
 \State Reject unless $a<b$
 \If{$a<0<b$} \State \Return $0$ \EndIf
 \If{$b\le0$} \State \Return $-\alg{POSITIVE}(-b,-a)$ \EndIf
 \State \Return $\alg{POSITIVE}(a,b)$
\ElsIf{$a=-\infty$ and $b=+\infty$} \State \Return $0$
\ElsIf{$a=-\infty$}
 \If{$b>0$} \State \Return $0$ \EndIf
 \State \Return $\lceil b\rceil-1$ \Comment{$\lceil b\rceil=-\lfloor-b\rfloor$}
\Else
 \If{$a<0$} \State \Return $0$ \EndIf
 \State \Return $\lfloor a\rfloor+1$
\EndIf
\end{algorithmic}
\end{samepage}

\begin{samepage}
\paragraph{Auxiliary procedure POSITIVE.}
Require finite $0\le a<b$. Store integer parts in a stack $K$; this implements
continued-fraction descent and recovery without recursive C++ calls.
\begin{algorithmic}[1]
\State $K\gets[\ ]$
\Loop
 \State $k\gets\lfloor a\rfloor$; $r\gets k+1$
 \If{$r<b$} \State \textbf{break} \EndIf
 \State Push $k$ onto $K$
 \State $f\gets a-k$; $a\gets1/(b-k)$
 \If{$f=0$}
  \State $r\gets\lfloor a\rfloor+1$; \textbf{break}
  \Comment{transformed upper endpoint is $+\infty$}
 \EndIf
 \State $b\gets1/f$
\EndLoop
\For{each $k$ in $K$ in reverse order} \State $r\gets k+1/r$ \EndFor
\State \Return $r$
\end{algorithmic}
\end{samepage}

\begin{lemma}[Exact simplest open-interval sample]\label{l:sampling}
For every nonempty open interval with rational finite endpoints and possibly
infinite outer endpoints, SAMPLE-RATIONAL terminates and returns its simplest
rational. In particular the result belongs strictly to the interval.
\end{lemma}
\begin{proof}
For POSITIVE, put $k=\lfloor a\rfloor$. If $k+1<b$, the returned integer
is strictly inside the interval. Otherwise $k\le a<b\le k+1$, and
$x\mapsto1/(x-k)$ is a bijection from $(a,b)$ onto
$(1/(b-k),1/(a-k))$, with inverse $y\mapsto k+1/y$.
If $a=k$, the latter upper endpoint is $+\infty$; the code immediately
chooses an integer there and never computes $1/0$.
The stack records the inverse maps, whose composition maps the current interval
onto the original positive interval. Reverse recovery therefore preserves strict
membership, including integer endpoints that must be excluded.

We prove the stronger statement that the positive result $p_0/q_0$ has both
$p_0\le p$ and $q_0\le q$ for every reduced positive candidate $p/q$.
In an integer base case, every candidate satisfies $p>kq$, hence
$p\ge kq+1\ge k+1$, and $q\ge1$.
For a nonterminal step, a reduced candidate $x=p/q$ transforms to the reduced
$y=q/(p-kq)$, with both parts positive. If the chosen transformed result
is $u/v$, induction gives $u\le q$ and $v\le p-kq$.
The recovered result is the reduced fraction $(ku+v)/u$:
its denominator is at most $q$, and its numerator is at most
$kq+(p-kq)=p$, since $k\ge0$.
Thus recovery preserves simultaneous minimality and in particular the specified
denominator-first order. Reflection preserves absolute numerators on negative
intervals. Zero is simplest whenever admitted. On an unbounded interval, the
specified integer is the admitted integer of least absolute value, so it too is
simplest.

For termination, write the finite current endpoints in reduced form as
$a=p/q$ and $b=s/t$. In a nonterminal step with noninteger $a$, the new
denominators are $s-kt$ and $p-kq$, with
$0<s-kt\le t$ and $0<p-kq<q$. Each is coprime to its corresponding
numerator, and their sum is strictly less than $q+t$.
An integer lower endpoint instead takes the immediate unbounded base case.
The positive integer measure $q+t$ therefore decreases until a base case is
reached. The finite stack is then recovered using positive divisors.
\end{proof}

For algebraic CAD boundaries, this lemma applies to the certified QQ inner gap,
not directly to the irrational endpoints. It guarantees simplestness in that
inner gap, not necessarily in the entire algebraic sector. For adjacent roots
$\alpha<\beta$, separation gives rational $a,b$ with
$\alpha<a<b<\beta$; hence $\alpha<a<q<b<\beta$.
Outer root brackets similarly give valid one-sided gaps. Section certificates
are unchanged. Sample selection imports no algebraic-coefficient operations.

\begin{algorithm}[H]
\caption{LIFT-STACK: one parent and its complete children}\label{a:stack}
\begin{algorithmic}[1]
\Require Parent sample $A$ of $S$ at depth $k-1$; current variable $x_k$;
ordered results $G=[(p_b,v_b)]_{b\in B_k}$ of LAZARD-QQ on $A$
\Ensure $2m+1$ ordered children, valid samples, dimensions, root sources and a stack record
\State $p\gets\prod_{b\in B_k}p_b$ in QQ polynomials \Comment{empty product is 1}
\State $R\gets\alg{ROOTS-QQ}(p,A,x_k)$ \Comment{sorted, distinct, complete root certificates}
\State Set per-source counters $N_b\gets0$
\For{each $D_j\in R$ in increasing order}
 \State $sources_j\gets\{b:\alg{POINT-ZERO}(p_b,D_j)\}$; require $sources_j\ne\varnothing$
 \State $b\gets\min(sources_j)$; $boundary_j\gets(b,N_b)$
 \For{$c\in sources_j$} \State $N_c\gets N_c+1$ \EndFor
\EndFor
\State $m\gets|R|$; $d\gets\dim(S)$ \Comment{$d=0$ for the depth-zero cell}
\For{$i=0,\ldots,m$}
 \If{$m=0$} \State $q\gets0$
 \ElsIf{$i=0$} \State $q\gets\alg{SAMPLE-RATIONAL}(-\infty,\inf(\text{last box of }D_1))$
 \ElsIf{$i=m$} \State $q\gets\alg{SAMPLE-RATIONAL}(\sup(\text{last box of }D_m),+\infty)$
 \Else
  \State Refine $D_i,D_{i+1}$ until their last boxes are strictly separated
  \State $q\gets\alg{SAMPLE-RATIONAL}(\sup(\text{last box of }D_i),\inf(\text{last box of }D_{i+1}))$
 \EndIf
 \State Append Sector $i$ to this parent's child block, with adjacent source/root boundary references
 \State Store sample $\alg{APPEND-RATIONAL}(A,x_k,q)$ and dimension $d+1$
 \If{$i<m$}
  \State Append Section $i+1$, sample $D_{i+1}$, reference $boundary_{i+1}$ and dimension $d$
 \EndIf
\EndFor
\State Store $R$, $G$ and $sources$ in the parent-indexed stack table
\end{algorithmic}
\end{algorithm}
APPEND-RATIONAL appends the QQ equation $x_k-q=0$, the strict interval
$[q-1,q+1]$, and a rational shift placing it in the positive half-line.
Its derivative is 1 and its unique root is $q$; existing prefix certificates are
retained. This proves validity and termination of this elementary helper.

\begin{algorithm}[H]
\caption{FULL-CAD: exact rational-data projection and lifting}\label{a:fullcad}
\begin{algorithmic}[1]
\Require Finite QQ input list $F$; projection variable order $V$
\Ensure A finite sign-invariant CAD of $\mathbb R^{|V|}$, its samples, dimensions and original-input leaf signs
\State Reject duplicate variables, unbound variables and negative exponents
\State Retain original QQ inputs; form denominator-cleared ZZ inputs $Z$
\State $P\gets\alg{PROJECT-FULL}(Z,V,\mathrm{LAZARD},\mathrm{FACTOR})$
\State Own the projection comparator; rebind returned polynomial comparators
\State Initialize the tree; reverse $V,P$ into lifting order $x_1,\ldots,x_n$ and $B_1,\ldots,B_n$
\For{$k=1,\ldots,n$}
 \State Parents are the preceding level in storage order, or the one virtual empty cell for $k=1$
 \For{each parent $S$ with sample $A$}
  \State $G\gets[\alg{LAZARD-QQ}(b,(x_1,\ldots,x_k),A):b\in B_k]$
  \State $\alg{LIFT-STACK}(S,A,x_k,G)$
 \EndFor
\EndFor
\For{each leaf, or the one empty leaf if $n=0$}
 \State Copy its sample to $A$; store $[\alg{POINT-SIGN}(f,A):f\in F]$
\EndFor
\State \Return the tree, samples, dimensions, stacks and signs
\end{algorithmic}
\end{algorithm}

\subsection{Loop invariants and implementation correspondence}
For LIFT-STACK, after $i$ iterations, the child block contains exactly the first
$i$ sectors and first $i$ sections (except that the final iteration adds only the
last sector). Every stored sample belongs to its child; boundary references preserve their
source polynomial index and per-source distinct-root ordinal; dimensions equal the parent dimension plus 0 or 1.
Refining a certificate changes neither its selected point nor its ordinal.
The source loop maintains $N_b$ equal to the number of roots already seen for
source $b$. Incrementing all sources, including secondary ones at shared roots,
proves that every stored reference has the baseline meaning.
An empty product has no roots and yields one sector, not an error.

For FULL-CAD, before level $k$, all prior levels partition their coordinate spaces,
have connected analytic semialgebraic cells and valid samples, and the immediately
preceding level is valuation-invariant for the nonzero members of $A_{k-1}$.
During a level, exactly the already visited parents have complete, contiguous
child blocks; each earlier child block stays unchanged as the next is appended.
Thus CSR parent/child indices, parallel sample and dimension arrays, and the stack
array indexed by parent remain aligned. The root-level virtual parent is encoded
by \texttt{SIZE\_MAX}, not by an ordinary node index. The sign loop has filled
exactly its visited leaf rows, each in original input order.

The storage type is \texttt{full\_cad\_result}, declared separately in
\texttt{full\_cad\_result.hh}. It only stores nodes and associated certificates;
its append operations keep the parallel tables aligned.
The corresponding algorithm code is \texttt{\_\_lift\_stack} and
\texttt{full\_cad}. Algebraic extensions use ROOTS-QQ; sector samples are chosen by
\texttt{sample\_rational} and certified by \texttt{append\_rational}.
Neither loop constructs an algebraic coefficient.
The Lazard result value type is declared separately in
\texttt{lazard\_qq\_result.hh}; its data consumer need not include the Lazard
evaluation algorithm header. The CAD result owns its comparator through a shared
pointer; projection polynomials, sample chains and root stacks borrow it. Copies
of the whole result retain this ownership. A detached point copy requires the
CAD comparator to remain alive. Lazard representatives separately retain shared
ownership of their own comparator.

\subsection{Correctness and termination}
\begin{theorem}[Complete rational-data CAD]\label{t:fullcad}
Relative to the stated trusted contracts for projection, factorization and tree storage, the
imported MPP theorem and the preceding real-point algorithms, FULL-CAD terminates
and returns a finite cylindrical algebraic decomposition covering all of
$\mathbb R^n$. Its cells are nonempty, connected, semialgebraic and pairwise
disjoint, its dimensions and samples are correct, and each recorded input sign
is constant over the corresponding leaf.
\end{theorem}
\begin{proof}
Start with $\mathbb R^0$, its single empty point and dimension 0. Its polynomial
data are constants and thus valuation-invariant. Suppose the driver invariant
holds before level $k$. The projection contract makes all required data for
$B_k$ valuation-invariant on each parent $S$. The product argument in
Section~\ref{s:cad} gives analytic ordered root functions and constant prefix
valuations on $S$. LAZARD-QQ correctly represents their specialization at the
sample using QQ polynomials. Specialization is a ring homomorphism, so the
product of representatives specializes to the product of their nonzero Lazard
values and is nonzero. ROOTS-QQ therefore applies and returns precisely their
$m$ distinct sample roots. This identifies all sections by order. Every such root
belongs to some factor's Lazard value, which proves that the computed source list
is nonempty. POINT-ZERO certifies each recorded source at that sample.

Sections and sectors partition every fibre over $S$: the $m$ equalities and
$m+1$ strict intervals exhaust the real line without intersection. Sections are
analytic graphs over $S$. Each sector is an open analytic submanifold of
$S\times\mathbb R$, homeomorphic to $S\times(0,1)$ using its continuous boundary
functions (with the usual one-sided or unbounded parametrization at infinity).
They are connected and nonempty, of dimensions $\dim S$ and $\dim S+1$,
respectively; their semialgebraicity was shown above. At the sample fibre the
strict rational brackets and Lemma~\ref{l:sampling} provide rational numbers
inside every sector, and the root certificates provide section samples.
APPEND-RATIONAL certifies the former.
Different parents are disjoint; each child's projection is exactly its parent.
Together these facts establish coverage, disjointness and cylindricality at the
new level as well as the child-block invariant.

By MPP Proposition 2.11 the product of $B_k$ is valuation-invariant on each child.
MPP Proposition 3.4 then gives valuation invariance of every individual factor on
each connected child. Lower-variable contents inherit their parent's valuation
with a trailing zero coordinate. Multiplicativity of valuation and the
factorization of each member of $A_k$ prove the driver's valuation invariant for
the new level. This also handles constants, omitted units and repeated factors.
The induction therefore continues through level $n$; it does not rely merely on
sign invariance of the parent samples.

A nonzero polynomial has valuation zero at a point exactly when its ordinary
value is nonzero. On a connected valuation-invariant cell it is consequently
identically zero, or nowhere zero. In the latter case continuity makes its sign
constant. Clearing positive denominators preserves the sign of each original
QQ input; zero inputs have sign zero everywhere. POINT-SIGN evaluates the original
input at its certified sample, and hence stores the sign of the whole leaf.
For $n=0$ this is rational constant evaluation at the sole empty leaf.

Projection is a terminating trusted operation. At each level there are finitely
many parents and projection factors. LAZARD-QQ and ROOTS-QQ terminate by their
proved contracts; each finite root list yields $2m+1$ children. Separation of two
distinct roots terminates by REFINE, and all source/sign queries terminate by
POINT-ZERO/SIGN. SAMPLE-RATIONAL terminates by Lemma~\ref{l:sampling}.
Induction over the finite number of levels and then the finite
leaf list proves termination of the driver and all its subloops.
\end{proof}
This theorem does not give a complexity bound or a machine-checked C++ refinement.
In particular, it imports the baseline algorithms' specified behavior rather than
claiming that wrapper checks prove their internals.

\section{Examples and mandatory regressions}
\begin{enumerate}
\item $T=[x^2-2]$, $B=[1,2]$, $f=(x^2-2)y$.
PREFIX-ZERO is true at order 0 and false for $2xy$.
Return representative $2xy$ and valuation $(1)$. The meaning is $2\sqrt2 y$;
no such coefficient is stored.
\item With the same $T,B$, $f=(x^2-2)^2$ gives $r=2$ and representative
$(12x^2-8)/2=6x^2-4$. Its value is $8$. Reducing $f$ modulo $T$ before
these derivatives would incorrectly erase the entire input.
\item $T=[x^2-2,z^2-2]$, positive $x$ and negative $z$.
POINT-ZERO on $z-x$ returns false, on $z+x$ returns true. Norm roots and
coordinate-interval overlap alone cannot distinguish the branches.
\item $f=xz-y$ at $(x,y)=(0,0)$, order $x,y,z$, returns $-1$ and $(0,1)$.
Do not specialize all rational coordinates in advance.
\item $T=[(x-1)(x-2)]$, a box around $1$, and $c=x-1$.
The representative and its quotient-ring residue may be nonzero although
POINT-ZERO must return true. No quotient-field assumption is permitted.
\item $T=[(x-1)(x-2)]$, selected $x=1$, $p=(x-2)y$.
ROOTS-QQ must remove the $x=2$ degree-zero branch before appending the equation.
Otherwise the whole line $x=2$ would violate the finite-solution hypothesis.
\item $T=[(x-1)^2]$ is regular but not squarefree. RCtoSqrfree must remove this
multiplicity even when the chain has length one.
\item Root isolation must include negative roots, zero roots, rational roots hit
by bisection, close roots, and specializations with degree drop or repeated roots.
\end{enumerate}

\section{Implementation audit and formal-proof boundary}
Review is against the numbered algorithms in this PDF, not against passing tests.
Each function must have a traceability entry: algorithm/line, concrete types,
preconditions, invariant, decreasing measure or convergence lemma, and regression.
A different backend is a failed conformance check until the specification and proof
are changed explicitly. No implicit change of mathematical representation is allowed.

\begin{longtable}{p{0.24\linewidth}p{0.31\linewidth}p{0.37\linewidth}}
\toprule Obligation & Implementation & Implementation contract\\\midrule
\endhead
Exact QQ intervals & \texttt{rational\_interval.hh} &
Independent arithmetic and set operations; no root or polynomial dependency.
Exact endpoint operations with explicit domain checks.\\
Triangular algorithms & \texttt{regular\_chain.hh}, backed by \texttt{charset.hh} &
XY16 WRSD and squarefree decomposition. The multiplicity scan includes
the first level. Chains borrow external comparators and validate contiguous prefix order.\\
XZ-LIFT / NSHR / REFINE & \texttt{real\_algebraic.hh} &
Rational bounds, positive shifts and strict root certificates retained;
bottom-up refinement follows the numbered pseudocode.\\
POINT-ZERO / SIGN & \texttt{point\_zero}, \texttt{point\_sign} &
WRSD set labels and interval elimination at the original selected point; no field
quotient or norm-root shortcut.\\
PREFIX-ZERO / LAZARD-QQ & \texttt{lazard\_qq.hh} &
Active-monomial grouping; frozen QQ representatives; ordered derivatives and
factorial division; no irrational-coefficient object.\\
REENCODE / ROOTS-QQ & \texttt{real\_algebraic.hh} &
Membership plus certified box comparison; WRSD degree branch selection before
squarefree isolation; exact duplicate removal.\\
Full lifting & \texttt{full\_cad.hh} &
Free lifting and driver functions use certified real points, trusted projection
and the storage interface of \texttt{full\_cad\_result}.\\
\bottomrule
\end{longtable}
For the first coordinate, the implementation invokes XZ-LIFT with an empty
prefix. Then $L=U$ is an exact rational polynomial, so this is UNI preceded by a
rational translation and followed by its inverse; it implements the first step
of ISOLATE-REGULAR with the same root set and additional explicit certificates.
REENCODE additionally checks point membership before box matching; this is a
sound strengthening of its precondition, not an interval-overlap shortcut.
Regression results are recorded separately in the current module test report.
This is regression evidence; inherited
triangular routines, UNI and the projection operator still need implementation
refinement proofs. The audit does not claim author-Maple cross-validation.


For formalization the dependency graph is acyclic:
\[
\begin{gathered}
\text{rational arithmetic, polynomials, interval inclusion}\\
\Downarrow\\
\text{XY16 decomposition and XZ06 certified refinement}\\
\Downarrow\\
\text{POINT-ZERO}\longrightarrow\text{PREFIX-ZERO}\longrightarrow\text{LAZARD-QQ}.
\end{gathered}
\]
\noindent\alg{ROOTS-QQ} additionally uses squarefree decomposition and NREALZERO.
The proofs above assume correctness of the explicitly cited algorithms, exact
rational primitives and univariate isolation. Formal verification must discharge
these imports and the C++ representation refinements. The MPP projection/delineability theorem is a named import;
Theorem~\ref{t:fullcad} proves lifting and the driver relative to it.
Machine refinement of the C++ implementation remains a separate obligation.
Current tests and manual conformance review do not establish machine-checked
correctness or large-instance performance guarantees.

\begin{thebibliography}{9}
\bibitem{MPP} S. McCallum, A. Parusi\'nski and L. Paunescu,
Validity proof of Lazard's method for CAD construction,
\emph{Journal of Symbolic Computation} 92 (2019), 52--69.
Version checked: \url{https://arxiv.org/abs/1607.00264v2},
Algorithm 2.12, Propositions 2.11 and 3.4, Theorem 5.1 and Section 5.3.
Already cited by the trusted baseline project research.

\bibitem{XY16} B. Xia and L. Yang, \emph{Automated Inequality Proving and
Discovering}, World Scientific, 2016. Chapters 2--3. Local source:
\texttt{docs/references/triangular-set/}, file beginning ``Xia [and] Yang - 2016''.
\bibitem{XZ06} B. Xia and T. Zhang, Real Solution Isolation Using Interval
Arithmetic, \emph{Computers and Mathematics with Applications} 52 (2006),
853--860. DOI: 10.1016/j.camwa.2006.06.003. Local source:
\texttt{Xia-Zhang-2006-Real-Solution-Isolation-Using-Interval-Arithmetic.pdf}.
\bibitem{XY02} B. Xia and L. Yang, An Algorithm for Isolating the Real Solutions
of Semi-algebraic Systems, \emph{Journal of Symbolic Computation} 34 (2002),
461--477. Supplementary reference for regularization; its old isolation procedure
is not the completeness basis used here.
\end{thebibliography}

\section{Checked interfaces and lexicographic implementation}
The public polynomial entry points check variable coverage, nonnegative exponents
and coordinate compatibility before conversion. Their internal \texttt{\_\_}
variants consume normalized QQ or ZZ polynomials in the specified lexicographic
order and use assertions for caller obligations. Data types retain representation,
ownership and storage invariants; decomposition, refinement and CAD construction
remain free algorithms.

For PREFIX-ZERO the implementation puts every active variable before the frozen
variables, preserving the latter's relative order. Equal active monomials are
therefore contiguous. Collecting each contiguous group produces exactly the same
coefficient polynomial as the abstract map in Algorithm~\ref{a:prefix}; deleting
the common active factor preserves its sorted unique frozen terms. No second
sort or map is needed. LAZARD-QQ calls this ordered core on its unreduced QQ
representative and uses the exact point-zero core on the complete coefficient.

XZ-LIFT uses ordered interval evaluation for its leading-coefficient and root-bound
calculations. After the two bound polynomials pass the squarefree test it calls
the asserted UNI core, avoiding duplicate release-mode gcd computations. These
changes preserve all acceptance predicates and refinement evidence above.

The CAD wrapper supplies the checked ZZ projection inputs and original QQ sign
inputs to the lexicographic driver. LIFT-STACK is an internal free algorithm.
Lazard representatives share an immutable owned comparator with the result,
and root-source tests use coefficients converted once to that order. The stored
root identities and cell semantics are unchanged.
\end{document}
